\title[Storage competes for the same surplus]{Storage vs.\ the Flexible Load:\\
The Referee's First Question, Answered}


\begin{abstract}
The first question a referee asks about a ``computing load absorbs renewable surplus and that
is good for carbon'' story is: what about storage? Storage absorbs the same surplus and
\emph{returns} it to displace fossil, whereas the load consumes it for computing and displaces
nothing on the grid. So the ``no load'' counterfactual over-credits the load. We concede the
point in full and then show it does not damage the paper's results, because it is aimed at a
claim the paper does not need. Three facts, all verified: (i) for pure surplus absorption
storage lowers grid carbon and the load does not---storage dominates the load \emph{as a
carbon technology}; (ii) the load's carbon-relevant object is not surplus absorption but
\emph{dispatch}, and the naive-versus-surplus-contract carbon gap is exactly invariant to how
much storage is installed; (iii) the load is the higher-\emph{value} user of a surplus MWh
precisely when computing value clears a sharp threshold $R>c+\eta\,[\phio c+(1-\phio)\bp]$,
which high-value compute meets and a low-value load does not. The load is not a carbon
mitigation technology; it is a computing demand whose dispatch is a carbon lever, and that
lever is orthogonal to storage.
\end{abstract}


\section{The model with both assets}
Keep the residual-demand primitive $x=D-\KR\tilde S$ with c.d.f.\ $F(\cdot;\KR)$. In a surplus
state ($x<0$) the curtailable surplus is $(-x)$. Two assets can take it. Storage of capacity
$\KS$ and round-trip efficiency $\eta\in(0,1)$ charges $\min((-x),\KS)$ and returns $\eta$ of
it in later peak states, \emph{displacing fossil}. The flexible load of capacity $\KM$
computes on surplus, earning the margin $R-c$ per MWh and \emph{displacing no grid fossil}.
Storage is the fossil-displacer, so we let it charge first; the load takes the residual
surplus $\big((-x)-\text{charge}\big)_+$. Grid carbon (fossil burned) is
\begin{equation}\label{st-eq:carbon}
  \E[\text{fossil}] \;=\; \underbrace{\E[\min(x_+,\KF)]}_{\text{consumer fossil}}
  \;+\;\underbrace{\E[\text{naive load burn}]}_{=0\text{ under surplus contract}}
  \;-\;\underbrace{\eta\,\E[\min((-x)_+,\KS)]}_{\text{storage displacement}} .
\end{equation}

\section{Storage dominates the load as a carbon technology}
\begin{proposition}[Per-MWh carbon: the referee is right]\label{st-prop:carbon}
A unit of surplus routed to storage lowers grid fossil by $\eta$ times the fraction it returns
into fossil-serving states; the same unit routed to the (surplus-contracted) load lowers grid
fossil by \emph{zero}. Storage strictly dominates the load in grid-carbon terms for surplus
use.
\end{proposition}
\begin{proof}
Immediate from \eqref{st-eq:carbon}: $\KS$ enters with a strictly negative sign, $\KM$ (surplus
contract, naive burn $=0$) does not enter at all. The load computes; it does not return energy
to the grid.
\end{proof}
Numerically, at $\KR=2,\KF=7,\eta=0.85$: installing $\KS=3$ cuts grid fossil from $3.9333$ to
$3.9207$ ($-0.0126$); installing a surplus-contracted $\KM=3$ leaves it at $3.9333$
($-0.0000$). \emph{The consequence is that the ``no load'' counterfactual of Week~3 overstates
the load's carbon benefit}: the honest benchmark credits storage with the fossil it would have
displaced, and against that benchmark the surplus load is carbon-neutral, not carbon-negative.

\section{The dispatch gap is storage-invariant}
The load's carbon-relevant object is not whether it exists but how it runs. Under the
\emph{naive} rule it burns fossil for computing (Week~2's floor); under the \emph{surplus
contract} it does not. That difference is what Week~3 and Week~4 are about, and it is
untouched by storage.
\begin{proposition}[The naive$-$surplus carbon gap does not depend on $\KS$]\label{st-prop:invariant}
Let $\Delta_{\mathrm{disp}}(\KS)=\E[\text{fossil}\mid\text{naive},\KS]-\E[\text{fossil}\mid\text{surplus},\KS]$.
Then $\Delta_{\mathrm{disp}}(\KS)=\E[\text{naive load burn}]$ for all $\KS$: the storage term in
\eqref{st-eq:carbon} is common to both dispatch modes and cancels.
\end{proposition}
\begin{proof}
Both modes share the same consumer fossil and the same storage displacement $-\eta\E[\min((-x)_+,\KS)]$;
they differ only in the naive burn term, which is present under naive dispatch and zero under
the surplus contract. Subtracting, the storage term cancels identically.
\end{proof}
Verified: at $\KS=0,2,4$ the gap is $+2.1291$ in every case, constant to the reported digits.

\begin{corollary}[The gap survives storage rivalry; the crowding is second order]\label{st-cor:rivalry}
Proposition~\ref{st-prop:invariant} lets storage charge first, so the surplus load never
cannibalizes storage. Grant instead the load \emph{co-priority}---it takes surplus up to $\KM$
and storage gets only the residual. Then the surplus contract does crowd out some storage
displacement, and the gap becomes
$\E[\text{naive burn}]-\eta\big(\E[\min((-x)_+,\KS)]-\E[\min(((-x)_+-\KM)_+,\KS)]\big)$.
The crowding term is \emph{second order}: the avoided naive burn accrues in \emph{fossil-serving}
states (where the naive load runs on fossil), while the crowded-out storage charge is lost in
\emph{surplus} states---disjoint events. Numerically the gap falls only from $+2.129$ to
$+2.116$ (a $0.6\%$ reduction) as $\KS$ ranges up to $10$: the crowded-out displacement is
$\eta\times0.013$, two orders of magnitude below the avoided burn $2.129$. So the dispatch
result is robust to the rivalry, not an artifact of the charging order.
\end{corollary}
So every carbon result the paper actually proves---the naive rule raises carbon
(Week~4), the surplus contract is its carbon-optimal dispatch (Week~3)---holds with storage in
the model, and holds whether or not the load competes with storage for the surplus. Storage
changes the \emph{level} of carbon, not the \emph{dispatch comparison}.

\section{Who should take the surplus: a value threshold}
Storage wins the carbon comparison; it does not follow that storage should take all the
surplus, because the two assets sell different products. Storage arbitrages energy; its
expected value per charged MWh is $\eta$ times the marginal displacement price---generation
cost $c$ in the common case, the outage price $\bp$ in the rare shortfall---so
$v_{\mathrm{stor}}=\eta\,[\phio_{\mathrm{dm}}c+(1-\phio_{\mathrm{dm}})\bp]$ using the
\emph{main-model} fossil cutoff $F(\KF)=\phio_{\mathrm{dm}}=(\bp-c-\CF)/\beta$ of the
double-margin note. (Using this cutoff, not the utility-operated $\phio$, is what makes the
value comparison consistent with the paper's actual model---an inconsistency worth flagging,
since it moves the verdict.) The load produces computing at margin $R-c$.
\begin{proposition}[The load is the higher-value surplus user iff computing clears a threshold]
\label{st-prop:threshold}
The flexible load is the higher-value user of a marginal surplus MWh iff
\[
  \boxed{\;R \;>\; c+\eta\big[\phio_{\mathrm{dm}}c+(1-\phio_{\mathrm{dm}})\bp\big]\;=:\;R^{\star}.\;}
\]
Below $R^{\star}$ storage should take the surplus; above it the load should. Both are used at
an interior joint optimum whenever the surplus exceeds either asset's capacity, so they are
\emph{complements sorted by the load's value}, not pure substitutes.
\end{proposition}
At the calibration ($c=0.5,\eta=0.85,\bp=4.7,\CF=0.6$, so $\phio_{\mathrm{dm}}=0.90$) the
threshold is $R^{\star}=1.28$, and $v_{\mathrm{stor}}=0.782$. The paper's own $R=1.4$
(hence $R-c=0.90>0.782$) \emph{clears} it: at the calibration the flexible load \emph{is} the
higher-value surplus user, and storage should take the surplus only for a lower-value load.
The verdict is margin-sensitive---plugging the utility-operated $\phio=0.857$ instead would
give $R^{\star}=1.44$ and reverse it---which is exactly why the value comparison must be run
in the main model. It also depends on how rare outage is:
\begin{center}\small
\begin{tabular}{cccc}
\toprule
$\Pr(\text{outage})=1-\phio$ & $v_{\mathrm{stor}}$ & vs $R-c=0.90$ & higher value\\
\midrule
$0.02$ & $0.50$ & $>$ & load\\
$0.05$ & $0.60$ & $>$ & load\\
$\mathbf{0.10}$ & $\mathbf{0.78}$ & $\mathbf{>}$ & \textbf{load $\leftarrow$ main model ($\phio_{\mathrm{dm}}=0.90$)}\\
$0.14$ & $0.93$ & $<$ & storage $\leftarrow$ utility-op ($\phio=0.857$)\\
$0.20$ & $1.14$ & $<$ & storage\\
\bottomrule
\end{tabular}
\end{center}
The main-model row ($\Pr(\text{outage})=1-\phio_{\mathrm{dm}}=0.10$) is the relevant one, and
there the load clears; the utility-operated row ($0.14$) is the mismatched comparison that
reverses it.

\section{What this settles}
\begin{remark}[The honest bottom line]
Storage is the better carbon technology; the flexible load is not a carbon technology at all.
Its carbon footprint in the model is a \emph{dispatch} property, and the paper's results are
dispatch results: the naive rule raises emissions, the surplus contract is the carbon-optimal
way to run a load that exists for computing, and that comparison is exactly invariant to
storage (Proposition~\ref{st-prop:invariant}). What the storage counterfactual correctly kills is
the weaker reading that a surplus-absorbing load is carbon-\emph{negative} versus doing
nothing---against the right benchmark (storage) it is carbon-neutral, and the paper should, and
now does, benchmark the surplus contract against the naive rule rather than against a vacuum.
The load earns its place on private value (Proposition~\ref{st-prop:threshold}), which high-value
compute clears; the grid should still prefer storage for the surplus when the load's computing
is cheap. This is a cleaner and more defensible claim than ``computing greens the grid,'' and
it is the one the model supports.
\end{remark}
\emph{Limits.} This is a static, expected-flow treatment of storage: it prices the returned
energy at the expected marginal displacement, not through an explicit intertemporal
state-of-charge process. A dynamic storage model with binding state-of-charge constraints
would change the \emph{magnitudes} $v_{\mathrm{stor}}$ and the displacement fraction, and could
introduce a scarcity value when storage is small; it would not change
Proposition~\ref{st-prop:invariant} (the dispatch gap cancels the storage term whatever its size)
nor the \emph{existence} of a value threshold in Proposition~\ref{st-prop:threshold}. The full
dynamic (renewable, fossil, storage, flexible-load) program is the natural next model and is
out of Role~1's scope.
